\documentclass[a4paper,11pt]{article}

\usepackage[utf8]{inputenc}
\usepackage[T1]{fontenc}
\usepackage{amsmath,amssymb}
\usepackage{booktabs}
\usepackage{geometry}
\usepackage{hyperref}
\usepackage{listings}
\usepackage{xcolor}
\usepackage{enumitem}

\geometry{margin=2.5cm}

\hypersetup{
    colorlinks=true,
    linkcolor=blue!60!black,
    urlcolor=blue!60!black,
    citecolor=blue!60!black,
}

\lstdefinestyle{python}{
    language=Python,
    basicstyle=\ttfamily\small,
    keywordstyle=\color{blue!70!black}\bfseries,
    stringstyle=\color{red!60!black},
    commentstyle=\color{green!50!black}\itshape,
    backgroundcolor=\color{gray!8},
    frame=single,
    framerule=0.4pt,
    rulecolor=\color{gray!40},
    numbers=left,
    numberstyle=\tiny\color{gray},
    numbersep=8pt,
    xleftmargin=1.5em,
    breaklines=true,
    showstringspaces=false,
    tabsize=4,
    aboveskip=1em,
    belowskip=1em,
}

\newcommand{\code}[1]{\texttt{#1}}
\newcommand{\mat}[1]{\mathbf{#1}}
\renewcommand{\vec}[1]{\boldsymbol{#1}}

\title{\textbf{pyLGC} --- Python Interface for LGC2\\[0.5em]
\large Least-Squares Geodetic Computation}
\author{Jürgen Gutekunst, Marco Carotta}
\date{\today}

\begin{document}
\maketitle
\tableofcontents
\newpage

%==========================================================================
\section{Introduction}
%==========================================================================

\code{pyLGC} is a Python module that provides access to the LGC2
least-squares adjustment engine. It loads a native shared library
(\code{pyLGC\_C.dll} on Windows, \code{libpyLGC\_C.so} on Linux) via
\code{ctypes}.

The module exposes the evaluator, which encapsulates the linearised
least-squares problem built from an LGC input file, and allows the user
to inspect, perturb, evaluate, and solve the underlying system.
For a detailed description of the underlying concepts and methods,
see~\cite{gutekunst_lgc2}.

This document is structured as follows:
Section~\ref{sec:ls} describes the underlying least-squares problem,
Section~\ref{sec:api} documents the Python API, and
Section~\ref{sec:workflow} provides examples of common workflows.

%==========================================================================
\section{The Least-Squares Problem}
\label{sec:ls}
%==========================================================================

This section summarises the nonlinear estimation problem solved by LGC2.
The formulation follows~\cite{barbier_ls}; see
also~\cite{gutekunst_lgc2} for a broader overview
and~\cite{wells_krakiwsky} for the general theory of least squares.

In the general theory of geodetic adjustment, the observation model can
take the implicit (combined) form $F(x, L+v) = 0$, where both
parameters and observations appear inside $F$. In LGC2, however, all
observation models are of \emph{parametric} form: $f(x) - (L+v) = 0$,
i.e.\ the model predicts the observations as a function of the
parameters alone. As a consequence, the matrix
$B = \partial F / \partial L$ is always $-I$ (the negative identity),
which simplifies the solution considerably.

%--------------------------------------------------------------------------
\subsection{The Nonlinear Estimation Problem}
%--------------------------------------------------------------------------

Given a parametric observation model $f\colon \mathbb{R}^{n_x} \to
\mathbb{R}^{n_L}$ that predicts measurement values from unknown
parameters, LGC2 solves the constrained nonlinear least-squares problem
\begin{align}
  \begin{split}
    & \min \;\tfrac{1}{2}\, v^T P\, v \\
    & \text{over } (x, v) \in \mathbb{R}^{n_x} \times \mathbb{R}^{n_L}, \\
    & \text{subject to} \\
    & \quad f(x) + B\,v - L = 0, \\
    & \quad C(x) = 0,
  \end{split}
  \label{eq:pe_problem}
\end{align}
where $x \in \mathbb{R}^{n_x}$ is the parameter vector,
$L \in \mathbb{R}^{n_L}$ the measurement vector,
$v \in \mathbb{R}^{n_L}$ the residuals,
$B \in \mathbb{R}^{n_L \times n_L}$ the observation coefficient matrix,
$C\colon \mathbb{R}^{n_x} \to \mathbb{R}^{n_C}$ the constraint function,
and $P \in \mathbb{R}^{n_L \times n_L}$ the weight matrix.
Since all observation models in LGC2 are of parametric form
$f(x) - (L + v) = 0$, the matrix $B$ is always $-I$ (negative identity).
The API nevertheless exposes $B$ and $B^{-1}$ for completeness.

%--------------------------------------------------------------------------
\subsection{Gauss--Newton Linearisation}
%--------------------------------------------------------------------------

Problem~\eqref{eq:pe_problem} is solved by the Gauss--Newton algorithm.
At iterate $x_k$, we define the misclosure vector as
\begin{equation}
    W_k = f(x_k) - L \in \mathbb{R}^{n_L},
    \label{eq:misclosure}
\end{equation}
which measures how well the current parameters explain the measurements.
Linearising $f$ around $x_k$ gives
\begin{equation}
    f(x_k + \Delta x) \approx f(x_k) + A_k\, \Delta x,
    \qquad
    A_k = \frac{\partial f}{\partial x}(x_k) \in \mathbb{R}^{n_L \times n_x}.
    \label{eq:Adef}
\end{equation}
Similarly, the constraint is linearised as
\begin{equation}
    C(x_k + \Delta x) \approx C(x_k) + A_{2,k}\, \Delta x,
    \qquad
    A_{2,k} = \frac{\partial C}{\partial x}(x_k)
        \in \mathbb{R}^{n_C \times n_x}.
\end{equation}
Substituting the linearised model into the observation equation
$f(x) + B\,v - L = 0$ gives
\begin{equation}
    f(x_k) + A_k\,\Delta x + B\,v - L \;\approx\; 0
    \qquad\Longrightarrow\qquad
    v \;\approx\; -B^{-1}\bigl(f(x_k) - L + A_k\,\Delta x\bigr)
      \;=\; -B^{-1}(W_k + A_k\,\Delta x).
\end{equation}
With $B = -I$ this simplifies to $v \approx W_k + A_k\,\Delta x$.
Inserting into the objective
$\min\,\tfrac{1}{2}\,v^T P\,v$ yields the linearised subproblem
\begin{align}
  \begin{split}
    \min \;&\tfrac{1}{2}\, (W_k + A_k\,\Delta x)^T P\,
        (W_k + A_k\,\Delta x)
    \quad \text{over } \Delta x, \\
    \text{s.t.}\quad
    & A_{2,k}\,\Delta x + C(x_k) = 0.
  \end{split}
  \label{eq:GN_subproblem}
\end{align}
This defines the matrices and vectors accessible through the
\code{pyLGC} API:
\begin{align}
    A &= \frac{\partial f}{\partial x}(x_k) \in \mathbb{R}^{n_L \times n_x},
        &\text{(design / Jacobian matrix)} \\
    B &= -I \in \mathbb{R}^{n_L \times n_L},
        &\text{(observation coefficient matrix)} \\
    W &= f(x_k) - L \in \mathbb{R}^{n_L},
        &\text{(misclosure vector)} \\
    A_2 &= \frac{\partial C}{\partial x}(x_k) \in \mathbb{R}^{n_C \times n_x},
        &\text{(constraint Jacobian)} \\
    W_2 &= C(x_k) \in \mathbb{R}^{n_C},
        &\text{(constraint misclosure)} \\
    P &\in \mathbb{R}^{n_L \times n_L}.
        &\text{(weight matrix)}
\end{align}

%--------------------------------------------------------------------------
\subsection{Solution of the Linearised Problem}
%--------------------------------------------------------------------------

Setting the gradient of the Lagrangian of~\eqref{eq:GN_subproblem} to
zero yields the KKT system
\begin{equation}
  \begin{pmatrix}
    A^T P\, A & A_2^T \\
    A_2       & 0
  \end{pmatrix}
  \begin{pmatrix}
    \Delta x \\ \lambda
  \end{pmatrix}
  = -
  \begin{pmatrix}
    A^T P\, W \\
    W_2
  \end{pmatrix},
  \label{eq:kkt}
\end{equation}
where $\lambda$ is the vector of Lagrange multipliers for the
constraints. In the unconstrained case ($n_C = 0$), this simplifies to
the normal equations
\begin{equation}
    A^T P\, A\;\Delta x = -A^T P\, W.
    \label{eq:normal}
\end{equation}
The parameters are then updated as $x_{k+1} = x_k + \Delta x$ and the
process repeats until convergence. The residuals at convergence are
$v = f(x^*) - L = W^*$.

%==========================================================================
\section{API Reference}
\label{sec:api}
%==========================================================================

%--------------------------------------------------------------------------
\subsection{Architecture}
\label{sec:arch}
%--------------------------------------------------------------------------

\code{pyLGC} is structured in two layers:

\begin{enumerate}
    \item \textbf{\code{pyLGC\_C.dll / libpyLGC\_C.so}} --- A shared
        library with a plain C~API (\code{extern "C"}) wrapping the C++
        LGC2 engine. It has no dependency on Python and can be called
        from any language with a C foreign-function interface.

    \item \textbf{\code{pyLGC.py}} --- A pure-Python module that loads
        the shared library via \code{ctypes} and exposes the
        \code{Evaluator} and \code{AdjustableObject} classes
        (\code{AdjustablePoint} and \code{AdjustableFrame} are retained
        as aliases for \code{AdjustableObject}). Arrays are returned as
        \code{numpy} arrays for efficient data transfer.
        Requires \code{numpy}; \code{scipy} is optional (for sparse
        matrix construction).
\end{enumerate}

Because the shared library has no compile-time dependency on Python
headers, it can be used from any language with a C foreign-function
interface.

%--------------------------------------------------------------------------
\subsection{Module-Level Types}
%--------------------------------------------------------------------------

\paragraph{\code{UEOIndices}} Named tuple holding the four key
dimensions of the problem:

\begin{center}
\begin{tabular}{cll}
    \toprule
    Symbol & Field & Meaning \\
    \midrule
    $n_x$ & \code{UIndex} & Number of unknowns (parameters) \\
    $n_L$ & \code{EIndex} & Number of equations (= number of observations) \\
    $n_L$ & \code{OIndex} & Number of observations \\
    $n_C$ & \code{CIndex} & Number of constraints \\
    \bottomrule
\end{tabular}
\end{center}

%--------------------------------------------------------------------------
\subsection{Class \code{Evaluator}}
%--------------------------------------------------------------------------

The central class. Each instance encapsulates one least-squares problem
loaded from an LGC input file.

\subsubsection*{Constructor}

\begin{description}[style=nextline, leftmargin=1.5em]
\item[\code{Evaluator(filePath: str)}]
    Load and parse the given LGC input file. Raises \code{RuntimeError}
    if the file cannot be opened or contains errors.
\end{description}

\subsubsection*{Problem Dimensions}

\begin{description}[style=nextline, leftmargin=1.5em]
\item[\code{getProblemDimensions() -> UEOIndices}]
    Return the dimensions $(n_x, n_L, n_L, n_C)$ of the problem. Available
    immediately after construction.
\end{description}

\subsubsection*{Parameter Management}

\begin{description}[style=nextline, leftmargin=1.5em]
\item[\code{getEstimatedParameters() -> numpy.ndarray}]
    Return the current parameter vector $x \in \mathbb{R}^{n_x}$ as a
    one-dimensional numpy array.

\item[\code{setParameters(para)}]
    Set the parameter vector. Accepts any sequence of length~$n_x$.
    Raises \code{RuntimeError} if the dimension is wrong.
    \emph{Note:} after calling \code{setParameters}, a new
    \code{evaluateAtParameters()} is required before the getters return valid
    results.
\end{description}

\subsubsection*{Evaluation}

\begin{description}[style=nextline, leftmargin=1.5em]
\item[\code{evaluateAtParameters() -> bool}]
    Linearise the functional model at the current parameter values.
    Computes all matrices and vectors. Returns \code{True} on success.
    Must be called before any matrix or vector getter.
\end{description}

\subsubsection*{Vectors}

\begin{description}[style=nextline, leftmargin=1.5em]
\item[\code{getMisclosure() -> numpy.ndarray}]
    Return the misclosure vector $W = f(x_k) - L \in \mathbb{R}^{n_L}$.
    Requires a preceding \code{evaluateAtParameters()}.

\item[\code{getConstraintMisclosure() -> numpy.ndarray}]
    Return the constraint misclosure $W_2 = C(x_k) \in \mathbb{R}^{n_C}$.
    Requires a preceding \code{evaluateAtParameters()}.
\end{description}

\subsubsection*{Sparse Matrices}

All sparse matrix getters return a tuple of five elements:
\begin{center}
    \code{(rows, cols, vals, nrows, ncols)}
\end{center}
where \code{rows}, \code{cols} are integer arrays, \code{vals} is a
float array (COO format), and \code{nrows}, \code{ncols} are the matrix
dimensions. See Section~\ref{sec:sparse} for how to construct
\code{scipy.sparse} matrices from these tuples.

\begin{description}[style=nextline, leftmargin=1.5em]
\item[\code{getFirstDesignMatrix()}]
    Design matrix (Jacobian) $A = \partial f / \partial x \in \mathbb{R}^{n_L \times n_x}$.

\item[\code{getSecondDesignMatrix()}]
    Observation coefficient matrix $B \in \mathbb{R}^{n_L \times n_L}$.
    Since LGC2 uses parametric models, $B = -I$ always (see
    Section~\ref{sec:ls}).

\item[\code{getConstraintDesignMatrix()}]
    Constraint Jacobian $A_2 = \partial C / \partial x \in \mathbb{R}^{n_C \times n_x}$.

\item[\code{getWeightMatrix()}]
    Weight matrix $P \in \mathbb{R}^{n_L \times n_L}$.
\end{description}

All require a preceding \code{evaluateAtParameters()}.

\subsubsection*{Solving}

\begin{description}[style=nextline, leftmargin=1.5em]
\item[\code{solve() -> (bool, numpy.ndarray)}]
    Run the full iterative least-squares solution. Returns a tuple
    \code{(ok, solution)} where \code{ok} indicates convergence and
    \code{solution} $\in \mathbb{R}^{n_x}$ is the estimated parameter
    vector. Can be called without a prior \code{evaluateAtParameters()}.
    \emph{Note:} after solving, \code{evaluateAtParameters()} must be called again
    before the matrix/vector getters are valid.
\end{description}

\subsubsection*{Observation Mapping}

\begin{description}[style=nextline, leftmargin=1.5em]
\item[\code{getObsIndexToLineNumber() -> dict}]
    Return a dictionary mapping internal observation indices to line
    numbers in the input file. Useful for tracing residuals back to
    specific observations.
\end{description}

\subsubsection*{Point and Frame Access}

\begin{description}[style=nextline, leftmargin=1.5em]
\item[\code{getPoint(name: str) -> AdjustableObject}]
    Look up an adjustable point by name. Raises \code{RuntimeError} if
    not found.

\item[\code{getFrame(name: str) -> AdjustableObject}]
    Look up an adjustable reference frame by name. Raises
    \code{RuntimeError} if not found.
\end{description}

The two methods perform distinct lookups in the underlying engine
(against \code{LGCAdjustablePoint} and
\code{TAdjustableHelmertTransformation} respectively), but both return
the same Python wrapper class (see Section~\ref{sec:adjobj}).

%--------------------------------------------------------------------------
\subsection{Class \code{AdjustableObject}}
\label{sec:adjobj}
%--------------------------------------------------------------------------

Represents an adjustable object in the network --- either a point or a
reference frame. Instances are obtained via \code{Evaluator.getPoint()}
or \code{Evaluator.getFrame()} and are valid as long as the parent
\code{Evaluator} exists.

\code{AdjustablePoint} and \code{AdjustableFrame} are retained as
aliases for \code{AdjustableObject}, so existing code that references
those names continues to work.

\begin{description}[style=nextline, leftmargin=1.5em]
\item[\code{getName() -> str}]
    Return the object's name.

\item[\code{getFirstUidx() -> int}]
    Return the index of this object's first unknown in the global
    parameter vector. Together with \code{getRelativeUnknIndices()},
    this identifies which elements of $x$ belong to this object.

\item[\code{getRelativeUnknIndices() -> numpy.ndarray}]
    Return the relative unknown indices (offsets from
    \code{getFirstUidx()}) as a 1-D integer array.
    For a fully free 3D point this is \code{[0, 1, 2]}; for a partially
    fixed point some components are absent; for a fully fixed object
    the array is empty.
    For an adjustable Helmert frame the indices select among up to
    seven parameters (3 translations, 3 rotations, 1 scale).

\item[\code{getEstVector() -> numpy.ndarray}]
    Return the current estimated parameter vector as a 1-D float array.
    For a 3D point this is always 3 elements regardless of which
    components are fixed; for a Helmert frame this is the current
    transformation parameter vector.
\end{description}

%--------------------------------------------------------------------------
\subsection{Error Handling}
\label{sec:errors}
%--------------------------------------------------------------------------

All errors are reported as Python \code{RuntimeError} exceptions with
descriptive messages. Common error conditions:

\begin{center}
\small
\begin{tabular}{ll}
    \toprule
    Situation & Example message \\
    \midrule
    File not found &
        \code{Cannot open file: network.lgc2} \\
    Malformed input file &
        \code{Errors found in the input file.} \\
    Wrong parameter dimension &
        \code{variable has wrong dimension, expected: 13, actual: 3} \\
    Getter before \code{evaluateAtParameters()} &
        \code{Must call evaluateAtParameters() before using getters} \\
    Unknown point name &
        \code{Point NOPE not found} \\
    Unknown frame name &
        \code{Frame NOPE not found} \\
    \bottomrule
\end{tabular}
\end{center}

%==========================================================================
\section{Typical Workflow and Examples}
\label{sec:workflow}
%==========================================================================

All code listings in this section are available as standalone Python
scripts in the \code{examples/} directory alongside this document.

The standard usage pattern follows a \emph{load -- set -- evaluate -- get}
cycle:

\begin{enumerate}[leftmargin=2em]
    \item \textbf{Load.} Create an \code{Evaluator} from an LGC input
        file. This parses the file, builds the internal data structures,
        and sets the provisional parameter values.
    \item \textbf{Set parameters} (optional). Replace the current
        parameter vector with custom values using
        \code{setParameters()}.
    \item \textbf{Evaluate.} Call \code{evaluateAtParameters()} to linearise the
        functional model at the current parameters. This computes the
        matrices $A$, $P$, $A_2$ and the
        misclosure vectors $W$, $W_2$.
    \item \textbf{Get results.} Retrieve any combination of matrices
        and vectors via the getter methods.
    \item \textbf{Repeat} steps 2--4 as needed (e.g.\ for iterative
        solving or sensitivity analysis).
\end{enumerate}

Alternatively, \code{solve()} performs the full iterative solution
internally, returning the converged parameter vector directly.

\lstinputlisting[style=python, caption={Basic workflow example (\code{examples/01\_basic\_workflow.py}).}]{examples/01_basic_workflow.py}

%--------------------------------------------------------------------------
\subsection{Sparse Matrix Usage with SciPy}
\label{sec:sparse}
%--------------------------------------------------------------------------

All sparse matrices are returned in COO (coordinate) format as plain
arrays, compatible with \code{scipy.sparse}:

\lstinputlisting[style=python, caption={Converting COO tuples to SciPy sparse matrices (\code{examples/02\_sparse\_matrices.py}).}]{examples/02_sparse_matrices.py}

%--------------------------------------------------------------------------
\subsection{Building the Normal Equations in SciPy}
\label{sec:normal}
%--------------------------------------------------------------------------

Given the matrices from the evaluator, the normal
equation~\eqref{eq:normal} can be assembled and solved in Python:

\lstinputlisting[style=python, caption={Forming and solving the normal equations (\code{examples/03\_normal\_equations.py}).}]{examples/03_normal_equations.py}

%--------------------------------------------------------------------------
\subsection{Gauss--Newton Iteration}
\label{sec:gn}
%--------------------------------------------------------------------------

The following example implements a simple Gauss--Newton solver using
the matrices provided by the evaluator. At each iteration the normal
equations~\eqref{eq:normal} are formed and solved, and the parameter
vector is updated until the step size falls below a tolerance.

\lstinputlisting[style=python, caption={Simple Gauss--Newton solver, validated against the built-in solver (\code{examples/04\_gauss\_newton.py}).}]{examples/04_gauss_newton.py}

%--------------------------------------------------------------------------
%==========================================================================
\begin{thebibliography}{9}

\bibitem{gutekunst_lgc2}
    J.~Gutekunst,
    \emph{LGC -- Position Estimation Software: Principal Concepts and Methods},
    CERN, 2024.
    \url{https://edms.cern.ch/document/3067818/1}

\bibitem{barbier_ls}
    M.~Barbier,
    \emph{Least Square Equations for LGC2},
    CERN, 2015.
    \url{https://edms.cern.ch/document/1475972/2}

\bibitem{wells_krakiwsky}
    D.~Wells and E.~Krakiwsky,
    \emph{The Method of Least Squares},
    University of New Brunswick, Department of Geodesy and Geomatics
    Engineering, Lecture Notes No.~18, 1971.

\end{thebibliography}

\end{document}
